\documentclass[12pt,a4paper]{article}

\usepackage[utf8]{inputenc}
\usepackage[T1]{fontenc}
\usepackage[spanish]{babel}
\usepackage{amsmath,amssymb}
\usepackage{geometry}
\usepackage{booktabs}
\usepackage{array}
\usepackage{parskip}
\usepackage{fancyhdr}
\usepackage{listings}
\usepackage{xcolor}
\usepackage{mathtools}

\geometry{margin=2.5cm, headheight=15pt}

\definecolor{codebg}{RGB}{245,245,245}
\definecolor{keywordcolor}{RGB}{0,100,180}
\definecolor{commentcolor}{RGB}{80,110,80}
\lstset{
  language=Python,
  backgroundcolor=\color{codebg},
  frame=single, rulecolor=\color{gray!40},
  basicstyle=\ttfamily\small,
  keywordstyle=\color{keywordcolor}\bfseries,
  commentstyle=\color{commentcolor}\itshape,
  numbers=left, numberstyle=\tiny\color{gray}, numbersep=8pt,
  breaklines=true, showstringspaces=false, tabsize=4,
}

\pagestyle{fancy}
\fancyhf{}
\rhead{Modelado y Simulaci\'on --- 22/04/2026}
\lhead{Punto 4 --- Docente: Omar C\'aceres}
\cfoot{\thepage}

\title{\textbf{Examen --- Modelado y Simulaci\'on}\\[0.3em]
       {\large Punto 4: Integraci\'on Doble por Monte Carlo}}
\author{Juan Bautista Espino}
\date{22 de abril de 2026}

\begin{document}
\maketitle\thispagestyle{fancy}

% ══════════════════════════════════════════════════════════════════════════════
\section*{Enunciado}
% ══════════════════════════════════════════════════════════════════════════════

\[
  I = \int_0^1 \int_0^2 e^{3(x-y)}\,dy\,dx
\]

\begin{enumerate}
  \item[(a)] Aproximar $I$ por Monte Carlo con \texttt{np.random.seed(0)}, $n=10\,000$.
        Resolver analíticamente y presentar tabla de estad\'isticos (media, varianza,
        desvici\'on est\'andar, error est\'andar, IC 95\%).
  \item[(b)] Demostrar que el error decrece como $1/\sqrt{n}$. Hallar la nueva muestra
        $n'$ para reducir el error a la mitad y verificar simulando.
\end{enumerate}

% ══════════════════════════════════════════════════════════════════════════════
\section*{Soluci\'on anal\'itica (verificaci\'on)}
% ══════════════════════════════════════════════════════════════════════════════

Se factoriza el integrando: $e^{3(x-y)} = e^{3x}\cdot e^{-3y}$.

\paragraph{Integral interior} (en $y$, con $x$ fijo):
\[
  \int_0^2 e^{3x}e^{-3y}\,dy
  = e^{3x}\!\left[-\frac{1}{3}e^{-3y}\right]_0^2
  = e^{3x}\cdot\frac{1}{3}\bigl(1 - e^{-6}\bigr)
\]

\paragraph{Integral exterior} (en $x$):
\[
  I = \frac{1-e^{-6}}{3}\int_0^1 e^{3x}\,dx
    = \frac{1-e^{-6}}{3}\cdot\left[\frac{e^{3x}}{3}\right]_0^1
    = \frac{(e^3-1)(1-e^{-6})}{9}
\]

Sustituyendo valores num\'ericos:
\[
  e^3 - 1 \approx 19.085\,537, \qquad
  1 - e^{-6} \approx 0.997\,521
\]

\[
  \boxed{I = \frac{(e^3-1)(1-e^{-6})}{9} \approx 2.115\,359}
\]

% ══════════════════════════════════════════════════════════════════════════════
\section*{Parte (a) --- Estimaci\'on por Monte Carlo ($n=10\,000$)}
% ══════════════════════════════════════════════════════════════════════════════

\subsection*{Fundamento te\'orico}

Para integrar $f(x,y)$ sobre el rect\'angulo $\mathcal{R} = [0,1]\times[0,2]$
de \'area $A = 1\times 2 = 2$, se generan $n$ puntos aleatorios uniformes
$(x_i, y_i)\sim\mathcal{U}(\mathcal{R})$ y se estima:

\[
  \hat{I} = A\cdot\frac{1}{n}\sum_{i=1}^{n} f(x_i, y_i)
           = A\cdot\bar{f}
\]

Por la Ley de Grandes N\'umeros, $\hat{I}\xrightarrow{n\to\infty} I$.

\subsection*{Implementaci\'on}

\begin{lstlisting}[caption={Monte Carlo con n = 10 000 (\texttt{seed=0})}]
import numpy as np

np.random.seed(0)
n = 10_000
A = 1.0 * 2.0   # area del rectangulo [0,1]x[0,2]

x = np.random.uniform(0, 1, n)
y = np.random.uniform(0, 2, n)
fxy = np.exp(3 * (x - y))

media_muestral = np.mean(fxy)
I_mc = A * media_muestral
varianza  = np.var(fxy, ddof=1)
desv_std  = np.std(fxy, ddof=1)
error_std = desv_std / np.sqrt(n)
IC_inf = I_mc - 1.96 * A * error_std
IC_sup = I_mc + 1.96 * A * error_std
\end{lstlisting}

\subsection*{Tabla de estad\'isticos}

\begin{center}
\renewcommand{\arraystretch}{1.4}
\begin{tabular}{lll}
\toprule
\textbf{Estad\'istico} & \textbf{S\'imbolo / F\'ormula} & \textbf{Valor} \\
\midrule
Media muestral & $\bar{f} = \frac{1}{n}\sum f(x_i,y_i)$ & $1.086\,293$ \\
Estimaci\'on MC & $\hat{I} = A\cdot\bar{f}$ & $2.172\,586$ \\
Varianza & $s^2 = \frac{1}{n-1}\sum(f_i-\bar{f})^2$ & $4.698\,309$ \\
Desv. est\'andar & $s$ & $2.167\,558$ \\
Error est\'andar & $\mathrm{SE} = s/\sqrt{n}$ & $0.021\,676$ \\
IC 95\% (inferior) & $\hat{I} - 1.96\cdot A\cdot\mathrm{SE}$ & $2.087\,618$ \\
IC 95\% (superior) & $\hat{I} + 1.96\cdot A\cdot\mathrm{SE}$ & $2.257\,555$ \\
Valor anal\'itico & $I = (e^3-1)(1-e^{-6})/9$ & $2.115\,359$ \\
Error absoluto & $|\hat{I} - I|$ & $0.057\,228$ \\
\bottomrule
\end{tabular}
\end{center}

El valor anal\'itico $I = 2.115\,359$ cae dentro del intervalo de confianza
$[2.088,\, 2.258]$. \checkmark

% ══════════════════════════════════════════════════════════════════════════════
\section*{Parte (b) --- Reducci\'on del error: $1/\!\sqrt{n}$}
% ══════════════════════════════════════════════════════════════════════════════

\subsection*{Demostraci\'on matem\'atica}

El error est\'andar del estimador Monte Carlo es:
\[
  \varepsilon(n) = A\cdot\frac{s}{\sqrt{n}}
\]

Supongamos que se desea un nuevo tama\~no $n'$ tal que el error se reduzca a
la mitad:
\[
  \varepsilon(n') = \frac{\varepsilon(n)}{2}
  \implies
  A\cdot\frac{s}{\sqrt{n'}} = \frac{1}{2}\cdot A\cdot\frac{s}{\sqrt{n}}
  \implies
  \sqrt{n'} = 2\sqrt{n}
  \implies
  \boxed{n' = 4n}
\]

Para reducir el error a la mitad, la muestra debe \textbf{cuadruplicarse}.
En general, para reducir el error en un factor $k$:
\[
  n' = k^2\cdot n
\]

Esto refleja la convergencia ``lenta'' de Monte Carlo: $\varepsilon \propto 1/\!\sqrt{n}$,
independientemente de la dimensi\'on del problema (ventaja sobre m\'etodos
determin\'isticos en dimensi\'on alta).

\subsection*{Verificaci\'on: simulaci\'on con $n' = 4n = 40\,000$}

\begin{lstlisting}[caption={Segunda simulación con n' = 40 000}]
np.random.seed(0)
n2 = 4 * n   # 40 000
x2 = np.random.uniform(0, 1, n2)
y2 = np.random.uniform(0, 2, n2)
fxy2   = np.exp(3 * (x2 - y2))
I_mc2  = A * np.mean(fxy2)
se2    = np.std(fxy2, ddof=1) / np.sqrt(n2)
error2 = A * se2
\end{lstlisting}

\begin{center}
\renewcommand{\arraystretch}{1.4}
\begin{tabular}{lcc}
\toprule
\textbf{Par\'ametro} & \textbf{$n = 10\,000$} & \textbf{$n' = 40\,000$} \\
\midrule
Estimaci\'on $\hat{I}$ & $2.172\,586$ & $2.107\,378$ \\
Error est\'andar       & $0.021\,676$ & $0.010\,477$ \\
$A\cdot\mathrm{SE}$    & $0.043\,351$ & $0.020\,955$ \\
IC 95\% & $[2.088,\; 2.258]$ & $[2.066,\; 2.148]$ \\
Error absoluto         & $0.057\,228$ & $0.007\,980$ \\
\midrule
Raz\'on de errores $\varepsilon_1/\varepsilon_2$ & \multicolumn{2}{c}{$2.069 \approx 2$ \checkmark} \\
\bottomrule
\end{tabular}
\end{center}

El error del estimador con $n'=40\,000$ es aproximadamente la mitad del obtenido
con $n=10\,000$ (raz\'on $\approx 2.07$), confirmando la relaci\'on te\'orica
$\varepsilon \propto 1/\!\sqrt{n}$.

\subsection*{Conclusi\'on}

\begin{enumerate}
  \item Monte Carlo estima $I \approx 2.173$ con $n=10\,000$; el valor exacto
        $I = 2.115$ queda dentro del IC 95\%.
  \item Para reducir el error a la mitad se necesitan $n' = 4n = 40\,000$ muestras.
  \item La simulaci\'on con $n'=40\,000$ confirma que el error se redujo de
        $\approx 0.043$ a $\approx 0.021$ (factor $\approx 2$).
  \item La ventaja de Monte Carlo sobre m\'etodos determin\'isticos crece con la
        dimensi\'on: el error es siempre $\mathcal{O}(1/\!\sqrt{n})$,
        sin importar la dimensi\'on del dominio de integraci\'on.
\end{enumerate}

\end{document}
